\documentclass[12pt]{article}
\usepackage[a3paper, margin=1in]{geometry}
\usepackage{amsmath, amssymb, amsthm, ulem}

\begin{document}

\section*{Domácí úkol VII}

Vypracoval: Daniel \textit{"Randál"} Ransdorf \hfill Podpis: \rule{4cm}{0.4pt}

\begin{enumerate}
  \item Označme jev $A_4$ "kód neobsahuje číslici 4", označme jev $A_7$
    "kód neobsahuje číslici 7". Elementárními jevy budou jednotlivé kódy,
    množinu všech elementárních jevů označme $X, |X|=10000$.
    Vypočítejme, kolik je kódů, které neobsahují
    alespoň jednu z těchto číslic pomocí principu inkluze/exkluze:

    \[
      |A_4\cup A_7| = |A_4| + |A_7| - |A_4\cap A_7|
        = 9^4 + 9^4 - 8^4 = 9026
    \]

    Zbytek kódů kritérium splňují. Vypočtěme tedy
    pravděpodobnost, že kódu chybí alespoň jedna z číslic:

    \[
      P(A_4\cup A_7) = \frac{|A_4\cup A_7|}{|X|} = \frac{9026}{10000}
    \]

    Dopočtěme pravděpodobnost, že se $A_4\cup A_7$ nestane, tedy že kód
    obsahuje číslici $4$ i číslici $7$:

    \[
      P(\neg (A_4\cup A_7)) = 1-P(A_4\cup A_7) = \frac{10000-9026}{10000}
      = \frac{974}{10000} = \underline{\underline{9.74\%}}
    \]

  \item Střední hodnota, kolik správných písemek odevzdá jednomu studentovi
    je $E[1] = \frac{1}{n}$, protože jen jedna z $n$ písemek studentovi patří.
    Pak pomocí linearity střední hodnoty získáme,
    že pro $n$ studentů je střední hodnota $E[n] = n\cdot E[1] = \underline{\underline{\;1\;}}$.

\end{enumerate}

\end{document}
